\documentclass[11pt,oneside]{book}

\usepackage[a4paper,margin=1in]{geometry}
\usepackage{amsmath,amssymb,amsthm,mathtools}
\usepackage{enumitem}
\usepackage{xcolor}
\usepackage{hyperref}
\usepackage{microtype}

\hypersetup{
  colorlinks=true,
  linkcolor=blue!55!black,
  urlcolor=blue!55!black,
  citecolor=blue!55!black
}

\emergencystretch=2em

\newtheorem{theorem}{Theorem}[chapter]
\newtheorem{proposition}[theorem]{Proposition}
\newtheorem{lemma}[theorem]{Lemma}
\newtheorem{corollary}[theorem]{Corollary}
\newtheorem{definition}[theorem]{Definition}
\newtheorem{remark}[theorem]{Remark}

\DeclareMathOperator{\Tr}{Tr}
\DeclareMathOperator{\diag}{diag}
\DeclareMathOperator{\rank}{rank}
\DeclareMathOperator{\card}{card}
\DeclareMathOperator{\dist}{dist}
\DeclareMathOperator{\Var}{Var}

\newcommand{\C}{\mathbb C}
\newcommand{\R}{\mathbb R}
\newcommand{\N}{\mathbb N}
\newcommand{\E}{\mathbb E}
\newcommand{\Prob}{\mathbb P}
\newcommand{\1}{\mathbf 1}
\newcommand{\eps}{\varepsilon}
\newcommand{\GammaPT}{\Gamma}
\newcommand{\HS}{\mathrm{HS}}
\newcommand{\op}{\mathrm{op}}
\newcommand{\dd}{\,\mathrm d}
\newcommand{\sphere}{\mathbb S}
\newcommand{\abs}[1]{\left|#1\right|}
\newcommand{\norm}[1]{\left\|#1\right\|}
\newcommand{\ip}[2]{\left\langle #1,#2\right\rangle}
\newcommand{\set}[1]{\left\{#1\right\}}
\newcommand{\given}{\,:\,}

\title{Appendix B: A Reader-First Mathematical Rewrite}
\author{A book-style translation of the formal Appendix B}
\date{\today}

\begin{document}
\frontmatter
\maketitle

\chapter*{Preface}

This manuscript rewrites Appendix B as ordinary mathematics.  The goal is to
give a self-contained proof that a reader can follow from the definitions,
without relying on hidden bookkeeping or unnamed background objects.

The central object is the normalized induced random state
\[
  \rho=XX^*,\qquad X\in M_{d^2,s}(\C),\qquad \norm{X}_{\HS}=1,
\]
and its partial transpose \(\rho^{\GammaPT}\).  For fixed \(k\) we study
\[
  f_k(X)=\Tr\bigl((\rho^{\GammaPT})^k\bigr).
\]
The main result is a local-Lipschitz concentration estimate at the natural
scale \(d^{-(2k-2)}\), using the full spherical dimension \(2d^2s\), which is
of order \(d^4\) when \(s\sim\lambda d^2\).

\tableofcontents

\mainmatter

\chapter{The Model and the Norms}

\section{Finite Hilbert spaces}

Let \(P,Q,\Sigma\) be finite sets.  In the balanced specialization one takes
\[
  P=Q=\{1,\ldots,d\},\qquad \Sigma=\{1,\ldots,s\}.
\]
Write
\[
  I=P\times Q,\qquad D=\card I,\qquad S=\card \Sigma .
\]
Thus \(D=d^2\) in the balanced case.  We use the complex Hilbert spaces
\(\C^I\) and \(\C^\Sigma\), both with their standard Hermitian inner products:
\[
  \ip{x}{y}_{\C^I}=\sum_{i\in I}x_i\overline{y_i},
  \qquad
  \ip{u}{v}_{\C^\Sigma}=\sum_{\alpha\in\Sigma}u_\alpha\overline{v_\alpha}.
\]

When an asymptotic estimate involves a parameter \(\lambda\), we specialize
further to a sequence \(s=s_d\) satisfying
\[
  \frac{s_d}{d^2}\longrightarrow \lambda\in(0,\infty).
\]
In that regime \(D=d^2\) and \(S=s_d\).  Constants written \(C_\lambda\) may
depend on \(\lambda\), but not on \(d\).

The sample matrix space is
\[
  \mathcal X=M_{I,\Sigma}(\C)\simeq \C^{I\times\Sigma}.
\]
We view \(G\in\mathcal X\) both as an array \((G_{i\alpha})\) and as a linear
map \(G:\C^\Sigma\to\C^I\).

\begin{definition}[Operator, Hilbert--Schmidt, and trace norms]
Let \(T:E\to F\) be a linear map between finite-dimensional complex Hilbert
spaces.

The operator norm is
\[
  \norm{T}_{\op}
  =
  \sup\{\norm{Tv}_2:\ v\in E,\ \norm v_2=1\}.
\]
Equivalently, \(\norm{T}_{\op}\) is the largest singular value of \(T\).

The Hilbert--Schmidt norm is
\[
  \norm{T}_{\HS}^2=\Tr(T^*T).
\]
In orthonormal bases this is
\[
  \norm{T}_{\HS}^2=\sum_{a,b}|T_{ab}|^2.
\]

The trace norm, or Schatten \(1\)-norm, is
\[
  \norm{T}_1=\Tr\bigl((T^*T)^{1/2}\bigr),
\]
the sum of the singular values of \(T\).
\end{definition}

For a sample matrix \(G\in\mathcal X\), this gives
\[
  \norm{G}_{\HS}^2
  =
  \sum_{i\in I}\sum_{\alpha\in\Sigma}|G_{i\alpha}|^2.
\]
For \(A\in M_I(\C)\), regarded as an operator on \(\C^I\), the same definitions
give \(\norm A_{\op}\), \(\norm A_{\HS}\), and \(\norm A_1\).  If \(A\) is
Hermitian with eigenvalues \(\lambda_1,\ldots,\lambda_D\), then
\[
  \norm A_{\op}=\max_j|\lambda_j|,
  \qquad
  \norm A_1=\sum_{j=1}^D|\lambda_j|,
  \qquad
  \norm A_{\HS}
  =
  \left(\sum_{j=1}^D|\lambda_j|^2\right)^{1/2}.
\]

\section{Partial transpose}

The index set \(I=P\times Q\) records the tensor decomposition
\(\C^I\simeq\C^P\otimes\C^Q\).  For \(A\in M_I(\C)\), its partial transpose
along the second tensor factor is the matrix \(A^{\GammaPT}\) defined by
\[
  A^{\GammaPT}_{(p,q),(p',q')}
  =
  A_{(p,q'),(p',q)}.
\]
Only the \(Q\)-coordinates are transposed.

\begin{lemma}[Basic norm facts]
For every \(A\in M_I(\C)\),
\[
  \norm{A^{\GammaPT}}_{\HS}=\norm A_{\HS},
  \qquad
  \norm A_{\op}\le \norm A_{\HS},
  \qquad
  \norm A_1\le \sqrt D\,\norm A_{\HS}.
\]
If \(A\) is Hermitian and \(m\) is even, then
\[
  \norm A_{\op}^m\le \Tr(A^m).
\]
\end{lemma}

\begin{proof}
The first identity holds because partial transpose only permutes matrix
coordinates.  The inequalities \(\norm A_{\op}\le\norm A_{\HS}\) and
\(\norm A_1\le\sqrt D\,\norm A_{\HS}\) are the elementary inequalities
\(\ell^\infty\le\ell^2\) and \(\ell^1\le\sqrt D\,\ell^2\) applied to the
singular values of \(A\).

If \(A\) is Hermitian, diagonalize it.  For even \(m\),
\[
  \Tr(A^m)=\sum_j\lambda_j^m=\sum_j|\lambda_j|^m
  \ge \max_j|\lambda_j|^m=\norm A_{\op}^m.
\]
\end{proof}

\chapter{Gaussian Vectors and Polar Coordinates}

\section{The standard complex Gaussian}

A standard complex Gaussian is a random variable
\[
  g=\frac{X+iY}{\sqrt2},
\]
where \(X,Y\) are independent real \(N(0,1)\) variables.  Then
\[
  \E g=0,\qquad \E |g|^2=1,
\]
and \(|g|^2\) has exponential distribution with mean \(1\).

Let \(G=(G_{i\alpha})\in\mathcal X\) have independent standard complex Gaussian
entries.  Its law is denoted by \(\gamma\).  The density of \(\gamma\) with
respect to Lebesgue measure on the real vector space underlying \(\mathcal X\)
is
\[
  \pi^{-DS}\exp(-\norm G_{\HS}^2).
\]
Consequently,
\[
  \norm G_{\HS}^2
  =
  \sum_{i\in I}\sum_{\alpha\in\Sigma}|G_{i\alpha}|^2
\]
has Gamma distribution with shape parameter \(DS\) and scale parameter \(1\).

For each \(i\in I\),
\[
  \sum_{\alpha\in\Sigma}|G_{i\alpha}|^2
\]
has Gamma distribution with shape parameter \(S\) and scale parameter \(1\).
This is the source of the diagonal Gamma estimates later.

\section{Wishart matrices and normalized states}

The unnormalized Wishart matrix is
\[
  GG^*\in M_I(\C).
\]
We also use the normalized Wishart matrix
\[
  W=\frac1S GG^*.
\]
The normalized induced state is obtained by radial normalization:
\[
  \rho(G)=
  \begin{cases}
    GG^*/\Tr(GG^*),&G\ne0,\\
    0,&G=0.
  \end{cases}
\]
Since \(\Tr(GG^*)=\norm G_{\HS}^2\), if
\[
  R(G)=\norm G_{\HS},
  \qquad
  \Theta(G)=\frac{G}{\norm G_{\HS}}\quad(G\ne0),
\]
then
\[
  \rho(G)=\Theta(G)\Theta(G)^*.
\]
The value of \(\Theta(0)\) may be chosen arbitrarily on the unit sphere,
because \(\gamma\{0\}=0\).

\section{The polar law}

Let
\[
  \sphere(\mathcal X)=\{X\in\mathcal X:\norm X_{\HS}=1\}
\]
and let \(\sigma_{\mathcal X}\) be the normalized surface measure on this
sphere.

\begin{theorem}[Gaussian polar law]
If \(G\) is a standard complex Gaussian vector in \(\mathcal X\), then
\[
  \Theta(G)=\frac{G}{\norm G_{\HS}}
\]
is distributed according to \(\sigma_{\mathcal X}\).  Moreover
\(R(G)=\norm G_{\HS}\) and \(\Theta(G)\) are independent.
\end{theorem}

\begin{proof}
The real dimension of \(\mathcal X\) is \(2DS\).  In polar coordinates
\[
  G=r\theta,\qquad r>0,\quad \theta\in\sphere(\mathcal X),
\]
Lebesgue measure factors as
\[
  \dd G=r^{2DS-1}\,\dd r\,\dd\sigma_{\mathcal X}(\theta)
\]
up to the normalizing surface-area constant.  The Gaussian density is radial:
\[
  \pi^{-DS}e^{-\norm G_{\HS}^2}
  =
  \pi^{-DS}e^{-r^2}.
\]
Thus the joint law of \((R,\Theta)\) has density proportional to
\[
  e^{-r^2}r^{2DS-1}\,\dd r\,\dd\sigma_{\mathcal X}(\theta).
\]
This is a product measure.  Hence \(R\) and \(\Theta\) are independent, and
the angular marginal is \(\sigma_{\mathcal X}\).
\end{proof}

\begin{remark}
Equivalently, the direction law is characterized as the unique probability
measure on the sphere invariant under the unitary group \(U(\mathcal X)\).
The Gaussian measure is unitary invariant, the normalization map
\(G\mapsto G/\norm G_{\HS}\) commutes with unitary transformations, and the
unitary group acts transitively on the sphere.  This gives the same
identification with surface measure.
\end{remark}

\begin{corollary}[Radial rescalings]
For \(G\ne0\),
\[
  GG^*=R(G)^2\Theta(G)\Theta(G)^*,
  \qquad
  (GG^*)^{\GammaPT}
  =
  R(G)^2(\Theta(G)\Theta(G)^*)^{\GammaPT}.
\]
Consequently,
\[
  \norm G_{\op}=R(G)\norm{\Theta(G)}_{\op},
\]
and
\[
  \norm{(GG^*)^{\GammaPT}}_{\op}
  =
  R(G)^2\norm{(\Theta(G)\Theta(G)^*)^{\GammaPT}}_{\op}.
\]
\end{corollary}

\chapter{Spherical Concentration and Localization}

\section{Medians}

Let \((\Omega,\mu)\) be a probability space and \(f:\Omega\to\R\) measurable.

\begin{definition}[Median]
A real number \(m\) is a median of \(f\) if
\[
  \mu\{f\le m\}\ge\frac12,
  \qquad
  \mu\{f\ge m\}\ge\frac12.
\]
\end{definition}

\begin{lemma}[Changing the center]
Let \(m\) be a median of \(f\).  For every \(a\in\R\) and every \(t>0\),
\[
  \mu\{|f-m|\ge t\}
  \le
  2\,\mu\{|f-a|\ge t/2\}.
\]
\end{lemma}

\begin{proof}
If \(|m-a|<t/2\), then \(|f-m|\ge t\) implies \(|f-a|\ge t/2\).  If
\(|m-a|\ge t/2\), then one of the two median half-spaces,
\(\{f\ge m\}\) or \(\{f\le m\}\), lies inside \(\{|f-a|\ge t/2\}\).  Hence
that latter event has probability at least \(1/2\), so twice its probability
is at least \(1\).
\end{proof}

\begin{lemma}[Median controlled by the mean]
If \(f\ge0\), then every median \(m\) of \(f\) satisfies
\[
  m\le 2\E f.
\]
\end{lemma}

\begin{proof}
By the median property, \(\mu\{f\ge m\}\ge1/2\).  Therefore
\[
  \E f\ge \int_{\{f\ge m\}} f\,\dd\mu
  \ge m\,\mu\{f\ge m\}
  \ge \frac m2.
\]
\end{proof}

\section{Levy's lemma on the Euclidean sphere}

Let \(\sphere^{N-1}\subset\R^N\) be the Euclidean unit sphere, equipped with
normalized surface measure \(\sigma_{N-1}\).

\begin{theorem}[Spherical Levy inequality]
Let \(g:\sphere^{N-1}\to\R\) be \(L\)-Lipschitz, with \(L>0\).  Then \(g\)
has a median \(M_g\) such that for every \(t>0\),
\[
  \sigma_{N-1}\{|g-M_g|\ge t\}
  \le
  2\exp\left(-\frac{Nt^2}{4L^2}\right).
\]
\end{theorem}

This is the only geometric concentration theorem needed.  The constants are
not optimized; what matters below is the dependence on \(Nt^2/L^2\).

\section{Localized Levy lemma}

\begin{theorem}[Localized Levy lemma]
Let \(\Omega_0\subset\sphere^{N-1}\) be measurable and let
\(\sigma_{N-1}(\Omega_0)\ge3/4\).  Let \(f:\sphere^{N-1}\to\R\) be measurable.
Assume that \(f\) is \(L\)-Lipschitz on \(\Omega_0\).  If \(m_f\) is a median
of \(f\), then for every \(t>0\),
\[
  \sigma_{N-1}\{|f-m_f|\ge t\}
  \le
  2\sigma_{N-1}(\Omega_0^c)
  +
  4\exp\left(-\frac{Nt^2}{16L^2}\right).
\]
\end{theorem}

\begin{proof}
Extend \(f|_{\Omega_0}\) to an \(L\)-Lipschitz function on the whole sphere by
McShane's formula
\[
  F(x)=\inf_{y\in\Omega_0}\bigl(f(y)+L\norm{x-y}_2\bigr).
\]
Then \(F=f\) on \(\Omega_0\).  If \(M_F\) is a median of \(F\), Levy's lemma
gives
\[
  \sigma_{N-1}\{|F-M_F|\ge u\}
  \le
  2\exp\left(-\frac{Nu^2}{4L^2}\right).
\]
Since \(F=f\) on \(\Omega_0\),
\[
  \sigma_{N-1}\{|f-M_F|\ge u\}
  \le
  \sigma_{N-1}(\Omega_0^c)
  +
  2\exp\left(-\frac{Nu^2}{4L^2}\right).
\]
Apply the center-changing lemma with \(a=M_F\) and \(u=t/2\):
\[
  \sigma_{N-1}\{|f-m_f|\ge t\}
  \le
  2\sigma_{N-1}\{|f-M_F|\ge t/2\}.
\]
This is the desired inequality.
\end{proof}

\section{From medians to means}

\begin{lemma}[Layer-cake mean estimate]
Assume that \(|f-m|\le R\) almost surely.  Let
\(\Phi:[0,R]\to[0,\infty)\) be Borel measurable with finite integral, and
assume that
\[
  \mu\{|f-m|\ge t\}\le \Phi(t)
  \qquad(0\le t\le R).
\]
Then
\[
  |\E f-m|
  \le
  \int_0^R\Phi(t)\,\dd t.
\]
\end{lemma}

\begin{proof}
The layer-cake identity gives
\[
  \E|f-m|
  =
  \int_0^\infty \mu\{|f-m|\ge t\}\,\dd t
  =
  \int_0^R \mu\{|f-m|\ge t\}\,\dd t.
\]
Finally, \(|\E f-m|\le\E|f-m|\).
\end{proof}

\chapter{Deterministic Lipschitz Calculus}

\section{Trace powers}

\begin{lemma}[Telescoping identity]
For \(A,B\in M_D(\C)\) and \(k\ge1\),
\[
  A^k-B^k
  =
  \sum_{j=0}^{k-1}A^j(A-B)B^{k-1-j}.
\]
\end{lemma}

\begin{proof}
Expand the sum.  Consecutive terms cancel, leaving \(A^k-B^k\).
\end{proof}

\begin{proposition}[Trace-power perturbation]
For \(A,B\in M_D(\C)\) and \(k\ge1\),
\[
  |\Tr(A^k)-\Tr(B^k)|
  \le
  k\max(\norm A_{\op},\norm B_{\op})^{k-1}\norm{A-B}_1.
\]
\end{proposition}

\begin{proof}
The nontrivial norm input is the finite-dimensional ideal property of the
trace norm:
\[
  \norm{RST}_1\le \norm R_{\op}\,\norm S_1\,\norm T_{\op}.
\]
Indeed, by duality between the trace norm and the operator norm,
\[
  \norm M_1=\sup_{\norm U_{\op}\le1}|\Tr(UM)|.
\]
Thus
\[
\begin{aligned}
  \norm{RST}_1
  &=
  \sup_{\norm U_{\op}\le1}|\Tr(URST)|  \\
  &=
  \sup_{\norm U_{\op}\le1}|\Tr(TURS)|  \\
  &\le
  \sup_{\norm U_{\op}\le1}\norm{TUR}_{\op}\,\norm S_1 \\
  &\le
  \norm T_{\op}\,\norm R_{\op}\,\norm S_1 .
\end{aligned}
\]
We also use \(\norm{A^j}_{\op}\le\norm A_{\op}^j\), by submultiplicativity of
the operator norm.

Now use the telescoping identity:
\[
  A^k-B^k=\sum_{j=0}^{k-1}M_j,
  \qquad
  M_j=A^j(A-B)B^{k-1-j}.
\]
By linearity of the trace, the triangle inequality, and
\(|\Tr M|\le\norm M_1\),
\[
\begin{aligned}
|\Tr(A^k)-\Tr(B^k)|
&=
\left|\sum_{j=0}^{k-1}\Tr(M_j)\right|\\
&\le
\sum_{j=0}^{k-1}|\Tr(M_j)|\\
&\le
\sum_{j=0}^{k-1}\norm{M_j}_1.
\end{aligned}
\]
Therefore
\[
\begin{aligned}
|\Tr(A^k)-\Tr(B^k)|
&\le
\sum_{j=0}^{k-1}\norm{A^j(A-B)B^{k-1-j}}_1\\
&\le
\sum_{j=0}^{k-1}
\norm A_{\op}^j\,\norm{A-B}_1\,\norm B_{\op}^{k-1-j}\\
&\le
k\max(\norm A_{\op},\norm B_{\op})^{k-1}\norm{A-B}_1.
\end{aligned}
\]
\end{proof}

\section{The partial-transpose moment map}

For \(X\in\sphere(\mathcal X)\), put
\[
  A_X=(XX^*)^{\GammaPT}.
\]
The function of interest is
\[
  f_k(X)=\Tr(A_X^k).
\]

\begin{proposition}[Deterministic Lipschitz estimate]
For \(X,Y\in\sphere(\mathcal X)\),
\[
\begin{aligned}
|f_k(X)-f_k(Y)|
&\le
k\sqrt D\,
\max(\norm{A_X}_{\op},\norm{A_Y}_{\op})^{k-1}\\
&\qquad\qquad
\times
(\norm X_{\op}+\norm Y_{\op})\,\norm{X-Y}_{\HS}.
\end{aligned}
\]
\end{proposition}

\begin{proof}
The trace-power perturbation gives
\[
  |f_k(X)-f_k(Y)|
  \le
  k\max(\norm{A_X}_{\op},\norm{A_Y}_{\op})^{k-1}
  \norm{A_X-A_Y}_1.
\]
Now
\[
  \norm{A_X-A_Y}_1
  \le
  \sqrt D\,\norm{A_X-A_Y}_{\HS}
  =
  \sqrt D\,\norm{XX^*-YY^*}_{\HS},
\]
because partial transpose is Hilbert--Schmidt isometric.  Finally
\[
  XX^*-YY^*
  =
  (X-Y)X^*+Y(X-Y)^*,
\]
so
\[
  \norm{XX^*-YY^*}_{\HS}
  \le
  (\norm X_{\op}+\norm Y_{\op})\norm{X-Y}_{\HS}.
\]
\end{proof}

\begin{corollary}[Lipschitz scale on the good set]
Assume \(D=d^2\).  On the set
\[
  \Omega(a,b)
  =
  \left\{
    X\in\sphere(\mathcal X):
    \norm X_{\op}\le\frac ad,\quad
    \norm{(XX^*)^{\GammaPT}}_{\op}\le\frac b{d^2}
  \right\},
\]
the function \(f_k\) is Lipschitz with constant
\[
  L_k\le 2kab^{k-1}\,d^{-(2k-2)}.
\]
\end{corollary}

\begin{proof}
Insert the two good-set bounds and \(\sqrt D=d\) into the deterministic
estimate.
\end{proof}

\chapter{The Wick Expansion and Aubrun Counting}

The off-diagonal estimate is the only genuinely combinatorial part of the
argument.  This chapter explains the objects being counted.  The point is not
that the formulas are pretty; the point is that after the Gaussian expansion,
everything reduces to counting how many choices remain after a finite list of
equality constraints has been imposed.

\section{Complex Wick formula}

Let \(z_\eta\), \(\eta\in\Xi\), be independent standard complex Gaussian
variables.

\begin{theorem}[Complex Wick formula]
For words \(h=(h_0,\ldots,h_{r-1})\) and
\(\bar h=(\bar h_0,\ldots,\bar h_{s-1})\) in \(\Xi\),
\[
  \E\prod_{a=0}^{r-1}z_{h_a}
    \prod_{b=0}^{s-1}\overline{z_{\bar h_b}}
  =
  0
\]
unless \(r=s\).  If \(r=s\), then the expectation is
\[
  \card\{\pi\in\mathfrak S_r:\ h_a=\bar h_{\pi(a)}
       \text{ for every }a\}.
\]
\end{theorem}

\begin{proof}
Group the factors by the coordinate \(\eta\).  For a fixed coordinate, a
nonzero expectation requires the same number of \(z_\eta\)'s and
\(\overline z_\eta\)'s, and the value is the factorial of that number.
Multiplying over coordinates counts precisely the bijective matchings between
holomorphic and anti-holomorphic positions that preserve the coordinate.
\end{proof}

\section{Closed walks and edge variables}

Let
\[
  Z=W^{\GammaPT}-\diag(W^{\GammaPT})
\]
be the off-diagonal part of the partially transposed Wishart matrix.  Since
\[
  W_{ij}=\frac1S\sum_{\alpha\in\Sigma}G_{i\alpha}\overline{G_{j\alpha}},
\]
the partially transposed entry is
\[
  (W^{\GammaPT})_{(p,q),(p',q')}
  =
  \frac1S\sum_{\alpha\in\Sigma}
  G_{(p,q'),\alpha}\,
  \overline{G_{(p',q),\alpha}}.
\]

Expanding \(\Tr(Z^m)\) gives a sum over closed walks
\[
  i_0,i_1,\ldots,i_m=i_0,
  \qquad
  i_r=(p_r,q_r)\in P\times Q,
\]
and over sample labels
\[
  \alpha_0,\ldots,\alpha_{m-1}\in\Sigma.
\]
For the edge from \(i_r\) to \(i_{r+1}\), the holomorphic Gaussian coordinate
is
\[
  \bigl((p_r,q_{r+1}),\alpha_r\bigr),
\]
and the anti-holomorphic coordinate is
\[
  \bigl((p_{r+1},q_r),\alpha_r\bigr).
\]
These crossed \(Q\)-coordinates are exactly where the partial transpose enters.
The off-diagonal projection means that the edge is discarded unless
\[
  i_r\ne i_{r+1}.
\]

\begin{definition}[Surviving Wick permutation]
Fix a closed walk \(i_0,\ldots,i_m=i_0\) and sample labels
\(\alpha_0,\ldots,\alpha_{m-1}\).  A permutation
\(\pi\in\mathfrak S_m\) is surviving if, for every edge \(r\),
\[
  (p_r,q_{r+1},\alpha_r)
  =
  (p_{\pi(r)+1},q_{\pi(r)},\alpha_{\pi(r)}),
\]
where the edge indices are read cyclically.  Equivalently,
\[
  p_r=p_{\pi(r)+1},\qquad
  q_{r+1}=q_{\pi(r)},\qquad
  \alpha_r=\alpha_{\pi(r)}.
\]
\end{definition}

The definition simply says that the holomorphic variable on edge \(r\) is
paired by Wick's formula with the anti-holomorphic variable on edge
\(\pi(r)\).

\begin{proposition}[Moment expansion after Wick's formula]
For every \(m\ge1\),
\[
  \E\Tr(Z^m)
\]
is the sum, over closed walks and sample labels satisfying the off-diagonal
condition, of the number of surviving Wick permutations, multiplied by the
normalizing factor \(S^{-m}\).
\end{proposition}

\begin{proof}
Expand the trace, expand \(W=S^{-1}GG^*\), and apply the complex Wick formula.
Each term in the expansion contains \(m\) Gaussian variables and \(m\)
conjugate Gaussian variables.  Wick's formula says that its expectation is the
number of permutations matching each holomorphic coordinate with an equal
anti-holomorphic coordinate.  Those are exactly the surviving permutations
defined above.
\end{proof}

\section{Relations generated by a surviving permutation}

Fix a Wick permutation \(\pi\).  Its equalities generate three equivalence
relations on the set of edge positions:
\[
  \sim_L,\qquad \sim_R,\qquad \sim_\Sigma.
\]
They are generated by
\[
  p_r=p_{\pi(r)+1},\qquad
  q_{r+1}=q_{\pi(r)},\qquad
  \alpha_r=\alpha_{\pi(r)}.
\]
Let
\[
  \ell_L(\pi),\qquad \ell_R(\pi),\qquad \ell_\Sigma(\pi)
\]
be the numbers of equivalence classes.  Once one value is chosen for each
class, all coordinates in that class are forced.  Therefore the number of
assignments compatible with \(\pi\) is at most
\[
  |P|^{\ell_L(\pi)}
  |Q|^{\ell_R(\pi)}
  |\Sigma|^{\ell_\Sigma(\pi)}.
\]
In the balanced case \(|P|=|Q|=d\), this becomes
\[
  d^{\ell_L(\pi)+\ell_R(\pi)}
  S^{\ell_\Sigma(\pi)}.
\]

\begin{definition}[Profile of a Wick permutation]
The profile of \(\pi\) is the pair
\[
  (u,v)
  =
  \bigl(\ell_L(\pi)+\ell_R(\pi),\,\ell_\Sigma(\pi)\bigr).
\]
Let \(N_m(u,v)\) be the number of Wick permutations of length \(m\) with
profile \((u,v)\).
\end{definition}

The whole combinatorial problem is now concentrated in the numbers
\(N_m(u,v)\).  The weighted relation sum has the form
\[
  \sum_{u,v} N_m(u,v)\,d^u\,S^v.
\]
The formal development proves the exact regrouping identity behind this
formula before applying Aubrun's bounds.

\section{Aubrun's tilde words and innovations}

Aubrun's combinatorial proof repackages the same equality constraints into
three words
\[
  \widetilde a,\qquad \widetilde b,\qquad \widetilde c,
\]
and two lists, traditionally called list (3) and list (7).  List (3) records
the labels directly seen along the walk.  List (7) records the labels after
the Wick permutation and the partial-transpose crossing have been applied.
In the notation above, these lists are just a different coordinate system for
the same \(P\)-, \(Q\)-, and \(\Sigma\)-labels.

There are three kinds of conditions.

\begin{enumerate}[label=(\roman*)]
\item The Wishart matching condition says that paired Gaussian variables have
matching sample labels.
\item The triple matching condition says that the left, right, and sample
labels match in the crossed pattern dictated by partial transpose.
\item The nonrepetition condition on \((\widetilde a,\widetilde b)\) encodes
the off-diagonal projection.
\end{enumerate}

An innovation is the first occurrence of a new equivalence class while the
lists are scanned from left to right.  If a position is not an innovation,
its label has already appeared and is forced by the previous class
representative.  A compatibility defect measures how many apparent choices
are lost when the \(A\)- and \(C\)-relations must both be compatible with the
same Wick relation.

\begin{lemma}[Aubrun Lemma 3.6]
Fix an innovation pattern with \(U\) innovations.  In the canonical model with
alphabet size at most \(k\), the number of compatible tilde-word assignments
is at most
\[
  k^{3U}.
\]
\end{lemma}

\begin{proof}
At each innovation one may choose at most one new value for each of the three
words \(\widetilde a,\widetilde b,\widetilde c\).  Each choice has at most
\(k\) possibilities.  Noninnovative positions are forced by the previously
chosen class representatives and by the matching conditions.  Therefore each
innovation contributes at most \(k^3\), giving \(k^{3U}\).
\end{proof}

\begin{lemma}[Fixed-defect class count]
For a fixed defect value \(\Delta\), the number of compatible relation classes
is bounded by
\[
  2^k(2k)^{36\Delta}.
\]
\end{lemma}

\begin{proof}
The innovation data are lifted from the edge positions to the vertex positions.
The compatibility defect controls the number of additional identifications
forced by the simultaneous \(A\)- and \(C\)-constraints.  Combining the
innovation count with the compatibility loss yields the displayed
fixed-defect bound.
\end{proof}

\begin{theorem}[Aubrun graduate relation counting]
There is a polynomial \(Q\), independent of \(d\) and \(S\), such that for
every \(m\ge1\) the weighted relation sum satisfies
\[
  \sum_{u,v}N_m(u,v)\,d^u\,S^v
  \le
  2^m(d+Q(m))^{m+2}(\sqrt S+Q(m))^m.
\]
\end{theorem}

\begin{proof}
Partition the Wick permutations by their profile and by their defect.  Lemma
3.6 controls the choices associated with innovations.  The fixed-defect class
count controls how many compatible relation classes occur at a given defect.
When the defect is small, the number of free classes is close to the maximal
one and the corresponding term is absorbed by the powers of \(d\) and
\(\sqrt S\).  When the defect is large, the factor \((2k)^{36\Delta}\) is
absorbed into the polynomial correction \(Q(m)\).  Summing over all profiles
gives the displayed bound.
\end{proof}

\chapter{Aubrun's Off-Diagonal Moment Estimate}

\section{The trace-moment envelope}

Let
\[
  W=\frac1S GG^*,\qquad
  Y=W^{\GammaPT},\qquad
  Z=Y-\diag(Y).
\]

\begin{theorem}[Off-diagonal moment bound]
There is an explicit polynomial \(Q\) such that for every \(m\ge1\),
\[
  \left|\E\Tr(Z^m)\right|
  \le
  \left(\frac2S\right)^m
  (d+Q(m))^{m+2}(\sqrt S+Q(m))^m
\]
in the balanced specialization \(D=d^2\).
\end{theorem}

\begin{proof}
Write a cyclic word in the bipartite index set \(I=P\times Q\) as
\[
  w=(i_0,i_1,\ldots,i_m),\qquad i_m=i_0,
\]
and write \(i_r=(p_r,q_r)\).  Expanding the trace gives
\[
  \Tr(Z^m)
  =
  \sum_{w}
  \prod_{r=0}^{m-1} Z_{i_r,i_{r+1}} .
\]
For an off-diagonal edge \(i_r\ne i_{r+1}\), the partial transpose of
\(W=S^{-1}GG^*\) has entry
\[
  Z_{i_r,i_{r+1}}
  =
  \frac1S
  \sum_{\alpha_r\in\Sigma}
  G_{(p_r,q_{r+1}),\alpha_r}\,
  \overline{G_{(p_{r+1},q_r),\alpha_r}}.
\]
For a diagonal edge \(i_r=i_{r+1}\), the corresponding entry of \(Z\) is
zero.  Thus, after expanding the \(m\) sample sums,
\[
  \E\Tr(Z^m)
  =
  S^{-m}
  \sum_w \mathbf 1_{\mathrm{off}}(w)
  \sum_{\alpha_0,\ldots,\alpha_{m-1}\in\Sigma}
  \E\prod_{r=0}^{m-1}
  G_{(p_r,q_{r+1}),\alpha_r}\,
  \overline{G_{(p_{r+1},q_r),\alpha_r}},
\]
where \(\mathbf 1_{\mathrm{off}}(w)\) is \(1\) when all edges are
off-diagonal and \(0\) otherwise.

Apply the complex Wick formula to the expectation in the last line.  A Wick
permutation \(\pi\in\mathfrak S_m\) survives precisely when, for every
edge \(r\),
\[
  (p_r,q_{r+1},\alpha_r)
  =
  (p_{\pi(r)+1},q_{\pi(r)},\alpha_{\pi(r)}),
\]
with the indices read cyclically.  Hence
\[
  \E\Tr(Z^m)
  =
  S^{-m}
  \sum_{\pi\in\mathfrak S_m}
  \sum_{(w,\alpha)\in\mathcal F_\pi}
  \mathbf 1_{\mathrm{off}}(w).
\]
Taking absolute values and using \(0\le\mathbf 1_{\mathrm{off}}\le1\),
\[
  \left|\E\Tr(Z^m)\right|
  \le
  S^{-m}\sum_{\pi\in\mathfrak S_m}\card\mathcal F_\pi.
\]

Now regroup the surviving pairs by profile.  If \(\pi\) has profile
\((u,v)\), then the fiber bound gives
\[
  \card\mathcal F_\pi\le d^uS^v,
\]
because one chooses one \(P\)- and one \(Q\)-coordinate for each free
left/right class and one sample label for each free sample class.  Therefore
\[
  \sum_{\pi\in\mathfrak S_m}\card\mathcal F_\pi
  \le
  \sum_{u,v}N_m(u,v)d^uS^v.
\]
The relation-counting theorem gives
\[
  \sum_{u,v}N_m(u,v)d^uS^v
  \le
  2^m(d+Q(m))^{m+2}(\sqrt S+Q(m))^m.
\]
Multiplying by \(S^{-m}\) gives
\[
  \left|\E\Tr(Z^m)\right|
  \le
  S^{-m}\,2^m(d+Q(m))^{m+2}(\sqrt S+Q(m))^m
  =
  \left(\frac2S\right)^m
  (d+Q(m))^{m+2}(\sqrt S+Q(m))^m.
\]
\end{proof}

\section{From moments to operator norm}

\begin{lemma}[Even-moment extraction]
If \(Z\) is Hermitian and \(m\) is even, then
\[
  \E\norm Z_{\op}
  \le
  \bigl(\E\Tr(Z^m)\bigr)^{1/m}.
\]
\end{lemma}

\begin{proof}
Pointwise, \(\norm Z_{\op}^m\le\Tr(Z^m)\).  Taking expectations and then the
\(m\)-th root gives the result.
\end{proof}

\begin{theorem}[Off-diagonal expectation bound]
Assume \(S/d^2\to\lambda>0\).  Then there is a constant \(C_\lambda\) such
that
\[
  \E\norm{W^{\GammaPT}-\diag(W^{\GammaPT})}_{\op}
  \le
  C_\lambda
\]
for all sufficiently large \(d\).
\end{theorem}

\begin{proof}
Choose even integers \(m=m(d)\) tending to infinity slowly enough that
\[
  Q(m)=o(d),
  \qquad
  \frac{\log d}{m}\to0.
\]
Applying the even-moment extraction to the moment envelope gives
\[
  \E\norm Z_{\op}
  \le
  \frac2S(d+Q(m))^{1+2/m}(\sqrt S+Q(m)).
\]
Since \(S\sim\lambda d^2\), the right side is bounded by a constant depending
only on \(\lambda\).
\end{proof}

\chapter{Diagonal Gamma Bounds and Wishart Expectations}

\section{Gamma tails}

The diagonal entries of \(W^{\GammaPT}\) are the same as those of \(W\):
\[
  (W^{\GammaPT})_{ii}
  =
  W_{ii}
  =
  \frac1S\sum_{\alpha\in\Sigma}|G_{i\alpha}|^2.
\]
Thus each \(W_{ii}\) is the average of \(S\) independent exponential random
variables of mean \(1\).

\begin{lemma}[Gamma tail for one average]
If
\[
  \xi=\frac1S\sum_{\alpha=1}^S E_\alpha,
\]
where the \(E_\alpha\) are independent exponential variables of mean \(1\),
then for every \(u\ge0\),
\[
  \Prob\left\{\xi\ge1+2\sqrt{\frac uS}+2\frac uS\right\}
  \le e^{-u}.
\]
\end{lemma}

\begin{proof}
Let
\[
  Y=S\xi=\sum_{\alpha=1}^S E_\alpha .
\]
For \(0\le\theta<1\),
\[
  \E e^{\theta Y}
  =
  \prod_{\alpha=1}^S \E e^{\theta E_\alpha}
  =
  (1-\theta)^{-S},
\]
because each \(E_\alpha\) has density \(e^{-x}\mathbf 1_{x\ge0}\).  Put
\[
  y=\sqrt{\frac uS}.
\]
If \(u=0\), the desired bound is just a probability bound by \(1\).  Assume
\(u>0\) and choose
\[
  \theta=\frac{y}{1+y}.
\]
Then \(0<\theta<1\).  Markov's inequality gives
\[
  \Prob\{\xi\ge 1+2y+2y^2\}
  =
  \Prob\{Y\ge S(1+2y+2y^2)\}
  \le
  \exp\{-S\theta(1+2y+2y^2)\}(1-\theta)^{-S}.
\]
Since \(1-\theta=(1+y)^{-1}\), the exponent is
\[
  -S\bigl(\theta(1+2y+2y^2)-\log(1+y)\bigr).
\]
Finally,
\[
  \theta(1+2y+2y^2)-y^2
  =
  y+\frac{y^3}{1+y}
  \ge y
  \ge \log(1+y),
\]
so
\[
  \theta(1+2y+2y^2)-\log(1+y)\ge y^2=\frac uS.
\]
Therefore the probability is at most \(e^{-u}\).
\end{proof}

\begin{lemma}[Maximum of the diagonal]
For every \(u\ge0\),
\[
  \Prob\left\{
    \max_{i\in I}W_{ii}
    \ge
    1+2\sqrt{\frac{u+\log D}{S}}
      +2\frac{u+\log D}{S}
  \right\}
  \le e^{-u}.
\]
\end{lemma}

\begin{proof}
For each fixed \(i\), the previous lemma with \(u+\log D\) gives
\[
  \Prob\left\{
    W_{ii}\ge
    1+2\sqrt{\frac{u+\log D}{S}}
      +2\frac{u+\log D}{S}
  \right\}
  \le e^{-u-\log D}=\frac{e^{-u}}{D}.
\]
The event for the maximum is contained in the union of these \(D\) events.
The union bound therefore gives \(D(e^{-u}/D)=e^{-u}\).
\end{proof}

\begin{theorem}[Diagonal expectation bound]
If \(D=d^2\) and \(S/d^2\to\lambda>0\), then
\[
  \E\norm{\diag(W^{\GammaPT})}_{\op}\le C_\lambda
\]
for all sufficiently large \(d\).
\end{theorem}

\begin{proof}
Let \(M=\norm{\diag(W^{\GammaPT})}_{\op}=\max_i W_{ii}\).  The deterministic
offset is
\[
  a=1+2\sqrt{\frac{\log D}{S}}+2\frac{\log D}{S}.
\]
Since \(\sqrt{u+\log D}\le\sqrt u+\sqrt{\log D}\), the maximum tail implies
\[
  \Prob\left\{
    M\ge a+2\sqrt{\frac uS}+2\frac uS
  \right\}\le e^{-u}
  \qquad (u\ge0).
\]
By the layer-cake formula applied to \((M-a)_+\),
\[
  \E(M-a)_+
  =
  \int_0^\infty \Prob\{(M-a)_+\ge y\}\,\dd y .
\]
With
\[
  h(u)=2\sqrt{\frac uS}+2\frac uS,
\]
the change of variables \(y=h(u)\) gives
\[
  \E(M-a)_+
  \le
  \int_0^\infty e^{-u}h'(u)\,\dd u
  =
  \frac{\sqrt\pi}{\sqrt S}+\frac2S.
\]
Thus
\[
  \E M
  \le
  1+2\sqrt{\frac{\log D}{S}}+2\frac{\log D}{S}
  +\frac{\sqrt\pi}{\sqrt S}+\frac2S.
\]
If \(D=d^2\) and \(S/d^2\to\lambda>0\), this upper bound is uniformly bounded
for all sufficiently large \(d\).
\end{proof}

\section{The full partially transposed Wishart expectation}

\begin{theorem}[Wishart-gamma expectation bound]
If \(D=d^2\) and \(S/d^2\to\lambda>0\), then
\[
  \E\norm{W^{\GammaPT}}_{\op}\le C_\lambda
\]
for all sufficiently large \(d\).
\end{theorem}

\begin{proof}
Use
\[
  W^{\GammaPT}
  =
  \diag(W^{\GammaPT})
  +
  \bigl(W^{\GammaPT}-\diag(W^{\GammaPT})\bigr)
\]
and the triangle inequality.  The diagonal term is bounded by the Gamma
estimate; the off-diagonal term is bounded by Aubrun's moment estimate.
\end{proof}

\chapter{Normalized Operator-Norm Inputs}

\section{Two Gaussian estimates}

We use two elementary estimates for the Gaussian matrix \(G\).  They are
included here because they are exactly what turns the Gaussian bounds into
normalized spherical bounds.

\begin{lemma}[Gaussian operator norm]
There is a universal constant \(C\) such that, for every
\(G\in M_{I,\Sigma}(\C)\) with independent standard complex Gaussian entries,
\[
  \E\norm G_{\op}\le C(\sqrt D+\sqrt S).
\]
\end{lemma}

\begin{proof}
We give a finite-net proof.  For unit vectors \(u\in\C^I\) and
\(v\in\C^\Sigma\), write
\[
  X_{u,v}
  =
  \sum_{i,\alpha}G_{i\alpha}v_\alpha\overline{u_i}
  =
  \ip{u}{Gv}.
\]
Since the entries of \(G\) are independent standard complex Gaussians,
\(X_{u,v}\) is a centered complex Gaussian with variance
\[
  \E|X_{u,v}|^2
  =
  \sum_{i,\alpha}|u_i|^2|v_\alpha|^2
  =
  \norm u_2^2\norm v_2^2
  =
  1.
\]
Thus \(|X_{u,v}|^2\) is exponential with mean \(1\), and for every \(r\ge0\),
\[
  \Prob\{|X_{u,v}|\ge r\}=e^{-r^2}.
\]

Let \(N_I\) be a \(1/4\)-net of the unit sphere of \(\C^I\), and let
\(N_\Sigma\) be a \(1/4\)-net of the unit sphere of \(\C^\Sigma\).  Viewing
\(\C^I\) and \(\C^\Sigma\) as real Euclidean spaces of dimensions \(2D\) and
\(2S\), the elementary volume argument gives
\[
  |N_I|\le 9^{2D},\qquad |N_\Sigma|\le 9^{2S}.
\]
Indeed, take a maximal \(1/4\)-separated subset; the balls of radius \(1/8\)
around its points are disjoint and lie in the ball of radius \(9/8\).

We next pass from the net to the full operator norm.  Let
\[
  M=\max_{u\in N_I,\ v\in N_\Sigma}|\ip{u}{Gv}|.
\]
For arbitrary unit \(u,v\), choose \(u_0\in N_I\), \(v_0\in N_\Sigma\) with
\(\norm{u-u_0}_2\le1/4\) and \(\norm{v-v_0}_2\le1/4\).  Then
\[
\begin{aligned}
  |\ip{u}{Gv}|
  &\le
  |\ip{u_0}{Gv_0}|
  +
  |\ip{u-u_0}{Gv}|
  +
  |\ip{u_0}{G(v-v_0)}|        \\
  &\le
  M+\frac14\norm G_{\op}+\frac14\norm G_{\op}.
\end{aligned}
\]
Taking the supremum over unit \(u,v\) gives
\[
  \norm G_{\op}\le 2M.
\]

By the union bound and the scalar tail above,
\[
  \Prob\{M\ge r\}
  \le
  |N_I||N_\Sigma|e^{-r^2}
  \le
  \exp\{2(D+S)\log 9-r^2\}.
\]
Set
\[
  a=\sqrt{2(D+S)\log 9}.
\]
Using \(\E M=\int_0^\infty\Prob\{M\ge r\}\,\dd r\),
\[
\begin{aligned}
  \E M
  &\le
  a+\int_a^\infty e^{a^2-r^2}\,\dd r       \\
  &=
  a+\int_0^\infty e^{-2at-t^2}\,\dd t
  \le
  a+\frac1{2a}.
\end{aligned}
\]
Since \(D,S\ge1\), the last term is bounded by a universal constant times
\(\sqrt D+\sqrt S\).  Finally,
\[
  \E\norm G_{\op}\le 2\E M\le C(\sqrt D+\sqrt S)
\]
for a universal constant \(C\).
\end{proof}

\begin{lemma}[Radius estimates]
For the standard Gaussian \(G\in\mathcal X\),
\[
  R(G)^2=\norm G_{\HS}^2
\]
has Gamma distribution with shape \(DS\) and scale \(1\).  In particular
\[
  \E R(G)^2=DS,
  \qquad
  \E R(G)\ge c\sqrt{DS}
\]
for a universal constant \(c>0\).
\end{lemma}

\begin{proof}
The Gamma law was proved from the coordinate description of the standard
complex Gaussian.  The identity \(\E R^2=DS\) is the mean of a Gamma variable
with shape \(DS\).  For the lower bound, use the Paley--Zygmund inequality on
\(R^2\), or equivalently the elementary Gamma estimate
\[
  \E\sqrt Z\ge c\sqrt n
  \qquad
  (Z\sim\Gamma(n,1)).
\]
The constant is universal because the ratio
\(\Gamma(n+1/2)/(\Gamma(n)\sqrt n)\) is bounded below for \(n\ge1\).
\end{proof}

\section{Expectation bounds on the sphere}

Let \(X\) be uniform on the Hilbert--Schmidt unit sphere
\(\sphere(\mathcal X)\).  Equivalently, one may start with a standard complex
Gaussian matrix \(G\) in \(\mathcal X\), discard the null event \(G=0\), and
set
\[
  R=\norm G_{\HS},\qquad \Theta=\frac{G}{R}.
\]
The polar law says that \(\Theta\) is uniform on \(\sphere(\mathcal X)\) and
independent of \(R\).  Thus \(X\) may be realized as \(\Theta\), and on this
realization \(G=RX\).

\begin{theorem}[First normalized expectation]
If \(D=d^2\) and \(S/d^2\to\lambda>0\), then
\[
  \E\norm X_{\op}\le \frac{C_\lambda}{d}.
\]
\end{theorem}

\begin{proof}
Using the polar realization just described, \(G=RX\) with \(R\) independent
of \(X\), and therefore
\[
  \E\norm G_{\op}=(\E R)\E\norm X_{\op}.
\]
The Gaussian operator-norm lemma gives
\[
  \E\norm G_{\op}\le C(\sqrt D+\sqrt S).
\]
Also \(R^2\) has Gamma distribution with shape \(DS\), so the radius estimate
gives
\[
  \E R\ge c\sqrt{DS}=c\,d\sqrt S.
\]
Therefore
\[
  \E\norm X_{\op}
  \le
  \frac{C(d+\sqrt S)}{c\,d\sqrt S}
  \le
  \frac{C_\lambda}{d}.
\]
\end{proof}

\begin{theorem}[Second normalized expectation]
If \(D=d^2\) and \(S/d^2\to\lambda>0\), then
\[
  \E\norm{(XX^*)^{\GammaPT}}_{\op}
  \le
  \frac{C_\lambda}{d^2}.
\]
\end{theorem}

\begin{proof}
The radial identity gives
\[
  (GG^*)^{\GammaPT}
  =
  R^2(XX^*)^{\GammaPT}.
\]
Since \(R\) and \(X\) are independent,
\[
  \E\norm{(GG^*)^{\GammaPT}}_{\op}
  =
  (\E R^2)\,
  \E\norm{(XX^*)^{\GammaPT}}_{\op}.
\]
Now \(\E R^2=DS=d^2S\), and
\[
  (GG^*)^{\GammaPT}=S W^{\GammaPT}.
\]
The Wishart-gamma expectation bound says
\[
  \E\norm{W^{\GammaPT}}_{\op}\le C_\lambda.
\]
Thus
\[
  \E\norm{(XX^*)^{\GammaPT}}_{\op}
  =
  \frac{S\,\E\norm{W^{\GammaPT}}_{\op}}{d^2S}
  \le
  \frac{C_\lambda}{d^2}.
\]
\end{proof}

\section{High-probability good sets}

\begin{theorem}[Good operator-norm sets]
There exist \(a,b,c>0\), depending only on \(\lambda\), such that for all
sufficiently large \(d\),
\[
  \sigma_{\mathcal X}
  \left\{X:\norm X_{\op}>\frac ad\right\}
  \le e^{-cd^2},
\]
and
\[
  \sigma_{\mathcal X}
  \left\{X:\norm{(XX^*)^{\GammaPT}}_{\op}>\frac b{d^2}\right\}
  \le e^{-cd^2}.
\]
\end{theorem}

\begin{proof}
For the first event, the map \(X\mapsto\norm X_{\op}\) is \(1\)-Lipschitz on
the Hilbert--Schmidt sphere.  Its median is at most twice its expectation, so
the expectation bound places the median at scale \(d^{-1}\).  Spherical Levy
concentration then gives an \(e^{-cd^2}\) tail, since the real dimension is
\(2DS\asymp d^4\) and the deviation is of size \(d^{-1}\).

For the second event, set
\[
  g(X)=\norm{(XX^*)^{\GammaPT}}_{\op}.
\]
The deterministic estimate
\[
  |g(X)-g(Y)|
  \le
  (\norm X_{\op}+\norm Y_{\op})\norm{X-Y}_{\HS}
\]
shows that \(g\) is \(2a/d\)-Lipschitz on the first good set.  The median of
\(g\) is at most twice its expectation, hence is at scale \(d^{-2}\).  Applying
the localized Levy lemma on the first good set gives the stated
\(e^{-cd^2}\) tail.
\end{proof}

\chapter{The Local Lipschitz Concentration Theorem}

\section{Statement}

Assume \(D=d^2\), \(S=s_d\), and
\[
  \frac{s_d}{d^2}\longrightarrow \lambda>0.
\]
Let \(X\) be uniform on the Hilbert--Schmidt unit sphere in
\(M_{d^2,s_d}(\C)\), and set
\[
  f_k(X)=\Tr\left(\left((XX^*)^{\GammaPT}\right)^k\right).
\]

\begin{theorem}[Local Lipschitz concentration]
Fix \(k\ge1\).  There are constants \(c,C>0\), depending only on \(k\) and
\(\lambda\), such that for all sufficiently large \(d\) and all
\(\eps>0\),
\[
  \Prob\left\{
    \abs{f_k(X)-\E f_k(X)}
    \ge
    \frac{\eps}{d^{2k-2}}
  \right\}
  \le
  C e^{-cd^2}
  +
  C\exp\left(-c\frac{d^4\eps^2}{k^2}\right).
\]
\end{theorem}

\section{Proof}

\begin{proof}
Let
\[
  \Omega
  =
  \left\{
    X:\norm X_{\op}\le \frac ad,\quad
    \norm{(XX^*)^{\GammaPT}}_{\op}\le \frac b{d^2}
  \right\}.
\]
The good-set theorem gives
\[
  \Prob(\Omega^c)\le C e^{-cd^2}.
\]
On \(\Omega\), the deterministic Lipschitz estimate gives
\[
  |f_k(X)-f_k(Y)|
  \le
  \frac{Ck}{d^{2k-2}}\norm{X-Y}_{\HS}.
\]
Thus \(f_k\) is locally Lipschitz on \(\Omega\) with constant
\[
  L_k=\frac{Ck}{d^{2k-2}}.
\]

The real dimension of the sphere is
\[
  N=2d^2s_d\asymp d^4.
\]
Apply the localized Levy lemma with
\[
  t=\frac u{d^{2k-2}}.
\]
For a median \(M_k\) of \(f_k\), this gives
\[
  \Prob\left\{
    |f_k(X)-M_k|\ge\frac u{d^{2k-2}}
  \right\}
  \le
  C e^{-cd^2}
  +
  C\exp\left(-c\frac{d^4u^2}{k^2}\right).
\]

It remains to replace the median by the mean.  First note that
\[
  \norm{(XX^*)^{\GammaPT}}_{\HS}
  =
  \norm{XX^*}_{\HS}
  \le
  \norm{XX^*}_1
  =
  \norm X_{\HS}^2
  =
  1.
\]
Thus the eigenvalues \(\alpha_j\) of \((XX^*)^{\GammaPT}\) satisfy
\[
  \sum_j|\alpha_j|^2\le1.
\]
For \(k=1\), \(f_1(X)=\Tr(XX^*)=1\).  For \(k\ge2\),
\[
  |f_k(X)|
  =
  \left|\sum_j\alpha_j^k\right|
  \le
  \sum_j|\alpha_j|^k
  \le
  \left(\sum_j|\alpha_j|^2\right)^{k/2}
  \le1.
\]
So \(|f_k|\le1\).

Integrating the median-centered tail gives
\[
  |\E f_k-M_k|
  \le
  C e^{-cd^2}
  +
  C\frac{k}{d^{2k}}.
\]
If this is at most \(\eps/(2d^{2k-2})\), then
\[
  \{|f_k-\E f_k|\ge \eps d^{-(2k-2)}\}
  \subset
  \{|f_k-M_k|\ge \tfrac12\eps d^{-(2k-2)}\},
\]
and the desired bound follows from the median tail.

If the mean-median gap is larger than \(\eps/(2d^{2k-2})\), then the displayed
bound on the gap forces
\[
  \eps\le Ckd^{-2}
\]
for all large \(d\).  In that case the exponential
\[
  \exp\left(-c\frac{d^4\eps^2}{k^2}\right)
\]
is bounded below by a positive constant, so after enlarging \(C\) the right
hand side of the desired estimate is at least \(1\).  The probability bound is
then trivial.
\end{proof}

\chapter{Final Form}

The preceding theorem is the normalized concentration result for partial
transpose moments.  It is stronger than the older unnormalized polynomial
concentration route in the regime relevant to the induced model.

\begin{theorem}[Appendix B concentration package]
Fix \(k\ge1\) and \(\lambda>0\), and let \(s_d/d^2\to\lambda\).  If
\[
  X\sim\sigma_{\mathcal X}
  \quad\text{on}\quad
  \{X\in M_{d^2,s_d}(\C):\norm X_{\HS}=1\},
\]
then, with
\[
  \rho=XX^*,
  \qquad
  f_k(X)=\Tr((\rho^{\GammaPT})^k),
\]
there are constants \(c,C>0\), depending only on \(k\) and \(\lambda\), such
that for all sufficiently large \(d\) and all \(\eps>0\),
\[
  \Prob\left\{
    \abs{\Tr((\rho^{\GammaPT})^k)
      -\E\Tr((\rho^{\GammaPT})^k)}
    \ge
    \eps\,d^{-(2k-2)}
  \right\}
  \le
  C e^{-cd^2}
  +
  C\exp\left(-c\frac{d^4\eps^2}{k^2}\right).
\]
\end{theorem}

\begin{proof}
This is exactly the local Lipschitz concentration theorem applied to
\(\rho=XX^*\).
\end{proof}

\backmatter

\chapter{Reader's Map}

The proof has five conceptual inputs.

\begin{enumerate}
\item The Gaussian polar law identifies the normalized Gaussian direction with
surface measure on the Hilbert--Schmidt sphere and proves independence from
the radius.
\item Spherical Levy concentration becomes useful only after localization,
because the moment map is not globally Lipschitz at the optimal scale.
\item The deterministic matrix calculus shows that the moment map is
Lipschitz on the set where both \(\norm X_{\op}\) and
\(\norm{(XX^*)^{\GammaPT}}_{\op}\) are controlled.
\item Aubrun's Wick-counting argument supplies the off-diagonal expectation
bound for the partially transposed Wishart matrix.
\item The diagonal Gamma estimate, polar factorization, and spherical
concentration turn the Gaussian/Wishart estimates into high-probability
operator-norm bounds on the normalized sphere.
\end{enumerate}

The chapters above present these inputs in reader order: first the model and
the geometric normalization, then the local concentration mechanism, then the
Wick-counting estimate, and finally the normalized concentration theorem.

\end{document}
